\documentclass[11pt]{article}

% =========================================================
%                    Page / Font / Typography
% =========================================================
\usepackage[T1]{fontenc}
\usepackage{lmodern}
\usepackage[a4paper,margin=1.15in]{geometry}
\usepackage{microtype}

% Optional highlighting (remove if unused)
\usepackage{soul}

% Global line spacing
\usepackage{setspace}
\setstretch{1.18}

% =========================================================
%                          Math
% =========================================================
\usepackage{mathtools}
\usepackage{amssymb,amsfonts,amsthm,amsopn}

% Display-math spacing
\setlength{\abovedisplayskip}{14pt}
\setlength{\belowdisplayskip}{14pt}
\setlength{\abovedisplayshortskip}{10pt}
\setlength{\belowdisplayshortskip}{10pt}

% Extra breathing room inside align/gather
\setlength{\jot}{8pt}

% Allow page breaks in long align blocks
\allowdisplaybreaks

% =========================================================
%                          Lists
% =========================================================
\usepackage{enumitem}

% =========================================================
%                    Theorems / Environments
% =========================================================
\usepackage{aliascnt}

\numberwithin{equation}{section}

\newtheorem{theorem}{Theorem}[section]

\newaliascnt{lemma}{theorem}
\newtheorem{lemma}[lemma]{Lemma}
\aliascntresetthe{lemma}

\newaliascnt{proposition}{theorem}
\newtheorem{proposition}[proposition]{Proposition}
\aliascntresetthe{proposition}

\newaliascnt{corollary}{theorem}
\newtheorem{corollary}[corollary]{Corollary}
\aliascntresetthe{corollary}

\theoremstyle{definition}

\newaliascnt{definition}{theorem}
\newtheorem{definition}[definition]{Definition}
\aliascntresetthe{definition}

\newaliascnt{example}{theorem}
\newtheorem{example}[example]{Example}
\aliascntresetthe{example}

\theoremstyle{remark}

\newaliascnt{remark}{theorem}
\newtheorem{remark}[remark]{Remark}
\aliascntresetthe{remark}
% =========================================================
%                    References
% =========================================================
\usepackage[hidelinks]{hyperref}
\usepackage[capitalize,noabbrev,nameinlink]{cleveref}

\crefname{theorem}{theorem}{theorems}
\Crefname{theorem}{Theorem}{Theorems}

\crefname{lemma}{lemma}{lemmas}
\Crefname{lemma}{Lemma}{Lemmas}

\crefname{proposition}{proposition}{propositions}
\Crefname{proposition}{Proposition}{Propositions}

\crefname{corollary}{corollary}{corollaries}
\Crefname{corollary}{Corollary}{Corollaries}

\crefname{definition}{definition}{definitions}
\Crefname{definition}{Definition}{Definitions}

\crefname{example}{example}{examples}
\Crefname{example}{Example}{Examples}

\crefname{remark}{remark}{remarks}
\Crefname{remark}{Remark}{Remarks}

\crefname{equation}{equation}{equation}

% =========================================================
%                       Headings / Layout
% =========================================================
\usepackage{titlesec}

% Make \paragraph a proper heading (not run-in)
\titleformat{\paragraph}[hang]{\bfseries}{\theparagraph}{0.75em}{}
\titlespacing*{\paragraph}{0pt}{0.75em}{0.35em}

% Optional: a bit more whitespace around section headers
\titlespacing*{\section}{0pt}{1.2em}{0.7em}
\titlespacing*{\subsection}{0pt}{1.0em}{0.5em}
\titlespacing*{\subsubsection}{0pt}{0.8em}{0.35em}

% =========================================================
%                       Notation helpers
% =========================================================
\DeclareMathOperator{\diver}{div}
\newcommand{\R}{\mathbb{R}}
\newcommand{\Ltwo}{L^2(\Omega)}
\newcommand{\Honez}{H_0^1(\Omega)}
\newcommand{\Hp}{\frac{\partial H}{\partial p}}
\newcommand{\norm}[1]{\left\lVert #1 \right\rVert}

% =========================================================
%                           Title
% =========================================================
\title{Existence of Solutions for a Stationary Mean Field Game System}
\author{Avi Wine}
\date{\today}

\begin{document}

% -------------------- Title + ToC -------------------------
\maketitle
\thispagestyle{empty}

\begin{abstract}
This report studies a stationary mean field game system given by a Hamilton--Jacobi--Bellman equation coupled with a Kolmogorov--Fokker--Planck equation. 
Under suitable assumptions, we establish the existence of a weak solution pair in $H_0^1(\Omega)\times H_0^1(\Omega)$, and we extend the argument to a variant with a nonlinear source term in the KFP equation.
\end{abstract}

\tableofcontents
\clearpage

% =========================================================
%                         CONTENT
% =========================================================

\section{Introduction}

\paragraph{Context and Motivation}
Mean field games is a branch of game theory devoted to the analysis of a special subset of games. These systems are distinguished by having a very large number of players, each of whom is considered almost insignificant, with minimal influence on the system as a whole. 
In this framework the discrete population of agents is replaced by a continuum distribution, allowing the system to be described using density functions and weak (distributional) solutions of partial differential equations. \\ \\
Furthermore, each player is a rational observer and the general trend around them informs their decision on how to act, with the aim of achieving an optimal outcome.
The famous example in economics of such a game, introduced by Aumann, is the behaviour of stock traders whom we treat as a continuum \cite{aumann}.
In this paper we consider the stationary equilibrium problem characterised by a coupled pair of equations: a Hamilton--Jacobi--Bellman (HJB) equation and a Kolmogorov--Fokker--Planck (KFP) equation.

\begin{equation}\label{eq:system}
\begin{aligned}
-\Delta u + H(x,\nabla u) + \lambda u &= f(x,m) && \text{in } \Omega \text{ HJB},\\
-\Delta m - \diver\!\bigl(m\,\Hp(x,\nabla u)\bigr) + \lambda m &= G(x) && \text{in } \Omega \text{ KFP}.
\end{aligned}
\end{equation}
This is supplemented with homogeneous Dirichlet boundary conditions $u=0$ and $m=0$ on $\partial\Omega$.

\paragraph{Interpretation of Equations}
Here $m(x)$ represents the player density and $u(x)$ the value function of a representative agent in the mean field game. In other words, $u(x)$ is a player's optimal solution achieved by minimising the expected cost at position $x$.
The term $\Hp$ denotes the partial derivative of $H$ with respect to the momentum variable (i.e.\ the $\nabla u$ argument).
HJB describes a representative player's optimal control behaviour for a given player density $m$. \\ \\
On the other hand, KFP gives the evolution of $m(x)$ when every player satisfies the optimal solution $u$.
Each player tries to minimise their cost function $f(x,m)$ (in HJB), and the optimal control for an individual player is given by $\alpha^\ast = -\Hp(x,\nabla u)$, which makes up the drift term in KFP.
The Laplacian terms are diffusion terms accounting for randomness (e.g.\ Brownian motion) and preventing perfect alignment of players even under a shared strategy.
The parameter $\lambda>0$ acts as a discount factor in the optimisation problem and plays an important analytical role by ensuring coercivity and boundedness of the associated operators.
Finally, $G(\cdot)$ is an exogenous term controlling entry/exit of players and therefore directly affects $m$ in KFP. \cite{lasry_lions, pierre}

\paragraph{Mathematical Contribution}
In this paper we analyse the coupled system \cref{eq:system} and investigate the existence of a weak solution pair $(u,m)\in\Honez\times\Honez$. 
Establishing the existence of a solution pair $(u,m)$ is important because such solutions correspond to Nash equilibria of the mean field game. 
The function $u$ determines the optimal strategy of a representative agent, while $m$ describes the distribution of agents when all players adopt this strategy. 
A solution of the coupled system therefore represents a self-consistent equilibrium of the game. 

While the general analytical framework relies on classical PDE tools such as the Lax--Milgram theorem, elliptic regularity, and Schaefer's fixed point theorem,
several intermediate estimates and calculations appearing in Sections~2.1--2.3 
are carried out explicitly for the present system. These derivations form an original technical contribution of this work.

\paragraph{Assumptions}
\begin{enumerate}[label=\arabic*., leftmargin=*, itemsep=0.35em]
  \item $\Omega$ is bounded and $\partial\Omega$ is smooth.
  \item $H(x,\nabla u):\Omega \times \R^n\to\R$ is
  \begin{itemize}[leftmargin=*, itemsep=0.25em]
    \item Lipschitz continuous with respect to $\nabla u$ with Lipschitz constant $\tilde{C}$,
    \item bounded by $C\bigl(|\nabla u|+1\bigr)$ for a constant $C$, for all $x\in \Omega$ and $\nabla u\in\R^n$.
  \end{itemize}
  \item The partial derivative $\Hp(x,\, \nabla u)$ is $L^2$ continuous and essentially bounded by some constant $L$.
  \item $\lambda >0$ is sufficiently large (in particular $\lambda>\max(\,\frac{L^2}{4},\,\frac{C^2}{2},\,\frac{\tilde{C}^2}{4})$).
  \item $f(x,m(x)):\Omega \times \Ltwo \to \Ltwo$ is
  \begin{itemize}[leftmargin=*, itemsep=0.25em]
    \item continuous with respect to $m$.
    \item bounded by $D\bigl(|m|+1\bigr)$ for a constant $D$, for all $x\in \Omega$ and $m\in \Ltwo$.
  \end{itemize}
  \item $G(x) \in \Ltwo$.
\end{enumerate}

\begin{example}
A Hamiltonian satisfying Assumption 2 is
\[
H(x,p) = \sqrt{1+|p|^2} + V(x) \quad \text{on } \Omega = B(0,1) \text{ in } \R^n,
\]
where $V \in L^\infty(\Omega)$. In this case, we get the associated drift 
\[\Hp = \frac{p}{\sqrt{1 + |p|^2}}.\]
\end{example}

\begin{example}
For a coupling term satisfying Assumption 5, one may consider
\[
f(x,m) = a(x)\,\phi(m),
\]
where $a(x) \in L^\infty(\Omega)$ and $\phi:\Ltwo\to\Ltwo$ is continuous.
\end{example}

\paragraph{Main Result}
\begin{theorem}[Existence of the MFG solution]\label{thm:main_result}
Let the Mean Field Game system be characterised by \cref{eq:system} and suppose the assumptions above hold.
Then a solution $(u,m) \in \Honez \times \Honez$ exists.
\end{theorem}

We prove this theorem throughout the rest of this paper.
\section{Method}

Our proof consists of two stages. First, we demonstrate solvability for each equation separately, while fixing the other variable (HJB with fixed $m$, and KFP with fixed $u$).
Second, we show the existence of a solution pair solving both equations simultaneously.
Before starting our proof it is important to introduce the following inequalities, which will be used repeatedly in this paper.

\begin{definition}
Parameters $p,q\in[1,\infty]$ are \emph{conjugate} if $\frac{1}{p}+\frac{1}{q}=1$.
\end{definition}

\begin{theorem}[Hölder's inequality]\label{thm:holder}
Let $f \in L^p(\Omega)$ and $g \in L^q(\Omega)$ where $p,q$ are conjugate. Then
\[
\norm{fg}_1 \le \norm{f}_p\,\norm{g}_q.
\]
\end{theorem}

\begin{remark}
In particular, when $p=q=2$ we obtain the Cauchy--Schwarz inequality
\[
\norm{fg}_1 \le \norm{f}_2\,\norm{g}_2.
\]
\end{remark}

\begin{theorem}[Young's inequality]\label{thm:youngs}
Let $a,b\in\R$. Then
\[
2ab \le \frac{a^2}{\varepsilon^2} + \varepsilon^2 b^2 \qquad \forall \varepsilon > 0.
\]
\end{theorem}

\subsection{Independent unique solutions}
\subsubsection{Solving the Kolmogorov--Fokker--Planck Equation}

To solve
\[
-\Delta m - \diver\!\bigl(m\,\Hp(x,\nabla u)\bigr) + \lambda m = G(x)
\]
with fixed $u$, we note that the equation is linear in $m(x)$, so we appeal to the Lax--Milgram theorem \cite{lax_milgram}.

\begin{theorem}[Lax--Milgram]\label{thm:lax-milgram}
Let $V$ be a Hilbert space and let $a:V\times V\to\R$ be a bilinear form such that for all $u,v\in V$,
\begin{enumerate}[label=(\roman*), leftmargin=*, itemsep=0.25em]
  \item (boundedness) $|a(u,v)|\le c_1\norm{u}\,\norm{v}$,
  \item (coercivity) $a(v,v)\ge c_0\norm{v}^2$.
\end{enumerate}
Let $l:V\to\R$ be a bounded linear functional. Then there exists a unique $u\in V$ such that
\[
a(u,v)=l(v)\qquad\text{for all }v\in V.
\]
Moreover, $\norm{u}_V \le \frac{1}{c_0}\norm{l}_{V^\ast}$.
\end{theorem}

In our case,
\begin{equation}\label{eq:weak_form}
a(m,v)=\int_\Omega \nabla m\cdot\nabla v + m\,\nabla v\cdot\Hp(x,\nabla u)+\lambda mv\,dx,
\qquad
l(v)=\int_\Omega G(x)\,v\,dx.
\end{equation}

We now show that KFP meets the conditions of \Cref{thm:lax-milgram}.

\begin{proposition}[Boundedness of KFP]\label{prop:kfp-bounded}
The form $a(m,v)$ is bounded on $\Honez\times\Honez$.
\end{proposition}

\begin{proof}
For arbitrary $m,v\in\Honez$,
\begin{align*}
|a(m,v)|
&\le \int_{\Omega} |\nabla m||\nabla v| +|\nabla v||m||\Hp| + \lambda |m||v| \,dx\\
&\le \norm{\nabla m}_{L^2}\norm{\nabla v}_{L^2} + L\norm{m}_{L^2}\norm{\nabla v}_{L^2} + \lambda\norm{m}_{L^2}\norm{v}_{L^2}.
\end{align*}
By Cauchy--Schwarz and Assumption 3. So there exists $c_1>0$ such that
\[
|a(m,v)| \le c_1\,\norm{m}_{H^1}\,\norm{v}_{H^1}.
\]
Thus $a(m,v)$ is bounded.
\end{proof}

\begin{proposition}[Coercivity of KFP]\label{prop:kfp-coercive}
If $\lambda > 0$ is sufficiently large, then $a(\cdot,\cdot)$ is coercive.
\end{proposition}

\begin{proof}
Consider
\begin{align*}
a(v,v)
&= \int_\Omega |\nabla v|^2 + v\,\nabla v\cdot\Hp(x,\nabla u)+\lambda v^2\,dx\\
&\ge \norm{\nabla v}_{L^2}^2 - L\norm{v}_{L^2}\norm{\nabla v}_{L^2} + \lambda\norm{v}_{L^2}^2.
\end{align*}
Applying Young's inequality yields
\begin{align*}
a(v,v)
&\ge \Bigl(1-\frac{\varepsilon^2}{2}\Bigr)\norm{\nabla v}_{L^2}^2
  + \Bigl(\lambda-\frac{L^2}{2\varepsilon^2}\Bigr)\norm{v}_{L^2}^2,
  \qquad \forall \varepsilon\in\R.
\end{align*}
To ensure $1-\varepsilon^2/2>0$, impose $\varepsilon^2<2$. Choosing $\varepsilon^2=2$ (as a limiting choice) gives the weakest requirement
\[
\lambda>\frac{L^2}{4},
\]
which ensures coercivity.
\end{proof}

The following theorem summarises the main result of this subsection. 
\begin{theorem}[Existence of unique KFP solution]\label{thm:solution_kfp}
Let the Kolmogorov--Fokker--Planck equation be as in \cref{eq:system} and suppose that $\lambda > 0$ is sufficiently large.
Then for any fixed $u \in \Honez$, there exists a unique $m \in \Honez$ solving the equation.
\end{theorem}

\begin{proof}
Let $a(m,v)=l(v)$ be the weak formulation as in \cref{eq:weak_form}. By \Cref{prop:kfp-bounded,prop:kfp-coercive}, $a(\cdot,\cdot)$ is bounded and coercive.
It remains to show $l$ is bounded. By Cauchy--Schwarz,
\begin{align*}
\norm{l}
&= \sup_{\norm{v}_{H^1}\le 1}\left|\int_\Omega Gv\,dx\right|
\le \sup_{\norm{v}_{L^2}\le 1}\norm{v}_{L^2}\,\norm{G}_{L^2}
\le \norm{G}_{\Ltwo}.
\end{align*}
Thus $l$ is bounded. By \Cref{thm:lax-milgram}, there exists a unique $m\in\Honez$ solving the KFP equation.
\end{proof}

\subsubsection{Solving Hamilton--Jacobi--Bellman}

Lax--Milgram does not apply directly to
\[-\Delta u + H(x,\nabla u) + \lambda u = f(x,m)\]
with fixed $m$, because the equation is nonlinear so we resort to a fixed point method instead. The idea is to construct a map from the HJB equation with the right properties such that 
it has a fixed point and for that very fixed point to be our solution. Remarkably, this fixed point method is general enough that we can 
use it to solve for existence in the system as whole. Indeed, this will be the concern of the next subsection. Throughout this section, $X$ denotes a real Banach space.
We now introduce the fixed point theorems that rely on suitable properties for the HJB context.

\begin{definition}
A mapping $A:X\to X$ is called \emph{compact} if for every bounded sequence $(u_k)_{k=1}^\infty$,
the sequence $(A[u_k])_{k=1}^\infty$ is precompact in $X$.
\end{definition}

\begin{theorem}[Schaefer's fixed point theorem]
  \label{thm:schaefer}
  Suppose $A:X\to X$ is continuous and compact, and assume the set
  \[\{u\in X: u=\mu A[u] \text{ for some } \mu \in [0,1]\}\]
  is bounded. Then $A$ has a fixed point. \footnote{See \cite[Section~9.2]{evans_pde}.}
\end{theorem}

Let us now design the map whose fixed point is a solution to the HJB equation. 
Consider

\[A:\Honez\to\Honez,\qquad A[u]=w \quad\text{where}\quad -\Delta w+\lambda w = f(x,m)-H(x,\nabla u).\]
We know $w$ exists by Lax--Milgram (\Cref{thm:lax-milgram}), so $A$ is well defined.\footnote{LHS is bounded and coercive  and RHS is bounded by linear growth condition on $H$.} Furthermore, by regularity theory \footnote{See \cite[Section~6.3]{evans_pde}.} we have 
\begin{align*}
  \norm{w}_{H^2(\Omega)}
  & \le C \norm{f(x,m)-H(x,\nabla u)}_{\Ltwo}\\
  & \le C \norm{f(x,m)}_{\Ltwo} + C\norm{H(x,\nabla u)}_{\Ltwo}\\
  & \le C \norm{f(x,m)}_{\Ltwo} + \hat{C}\norm{|\nabla u| + 1}_{\Ltwo} \qquad \text{(Assumption 2)}
\end{align*}
Where $C$, $\hat{C}$ are constants - note that these constants are not the same as the 
ones appearing in our assumptions above.
We now prove the existence of an HJB solution by way of our fixed point map.
\begin{theorem}[Existence of HJB Solution]
  \label{thm:existence_hjb}  
  If $\lambda>0$ is sufficiently large, then for any fixed $m \in \Honez$ there exists $u\in H^2(\Omega)\cap \Honez$ solving HJB.
\end{theorem}
\begin{proof}
  \footnote{Proof is adapted from \cite[Section~9.2]{evans_pde}, with several modifications specific to the present setting.}
  In other words, for sufficiently large $\lambda >0$, $A$ has a fixed point. 
  We show first that $A$ meets the right criteria for Schaefer's Theorem. We assert that $A$ is continuous and
  compact.\\
  If $u_k \to u$ in $\Honez$ and we denote $w_k = A[u_k]$ then higher regularity tells us 
  \[\sup_k \norm{w_k}_{H^2(\Omega)} < \infty\]
  (i.e $(w_k)$ is a bounded sequence in $H^2(\Omega) \cap \Honez$) and there exists a convergent subsequence $(w_{k_j})_{j=1}^{\infty} \to w \text{ in } \Honez$ \footnote{See \Cref{cor:embed}}.\\
  Now the weak form of the HJB equation with $w_{k_j}$ is
  \begin{align*}
    \int_\Omega \nabla w_{k_j}\cdot \nabla v + \lambda w_{k_j} \cdot v \,dx = \int_\Omega (f(x,m) - H(x,\nabla u_{k_j}))\cdot v \,dx, \text{ for any } v \in \Honez\\
  \end{align*}
  and consequently
  \begin{align*}
    \int_\Omega \nabla w\cdot \nabla v + \lambda w \cdot v \,dx = \int_\Omega (f(x,m) - H(x,\nabla u))\cdot v \,dx, \text{ for any } v \in \Honez\\
  \end{align*}
  by Lipschitz continuity of $H$ (Assumption 2).
  Thus, $w = A[u]$. Hence, $A[u_k] \to A[u]$ in $\Honez$, so $A$ is continuous. $A$ is compact by similar reasoning since 
  a bounded sequence $(u_k)_{k=1}^\infty$ yields a bounded sequence $(A[u_k])_{k=1}^\infty$ in $H^2(\Omega) \cap \Honez$ which has a strongly convergent subsequence in $\Honez$.

  It remains to show
  \[\{u\in \Honez: u=\mu A[u]\ \text{for some } \mu \in [0,1]\}\]
  is bounded.

  Substituting $w =u/\mu$ \footnote{\label{fn:mu} We ignore $\mu =0$ since then $w=0$.} and setting $v = u$, we obtain
  \begin{align*}
    \int_\Omega |\nabla u|^2 + \lambda |u|^2\,dx
    &= \mu\int_\Omega f u - H(x,\nabla u)u\,dx\\
    &\le \int_\Omega |fu| + |H(x,\nabla u)u|\,dx, \qquad \text{ since } 0< \mu \le 1\\
    &\le \int_\Omega |fu| + C(|\nabla u|+1)|u|\,dx \qquad \text{(Assumption 2)}\\
  \end{align*}
  To isolate $|\nabla u|^2$ and $u^2$ terms on the RHS, we apply Young's inequality (\Cref{thm:youngs}) once on $|fu|$ and twice over on $C(|\nabla u|+1)|u|$, to obtain
  \begin{align*}
    |fu|
    &\le \frac12(\frac{f^2}{\varepsilon^2} + \varepsilon^2 u^2) \qquad \varepsilon\in (0, \infty)
  \end{align*}
  and
  \begin{align*}
    C(|\nabla u|+1)|u|
    &\le \frac12\,\bigl(\frac{1}{\delta^2}\,|Cu|^2 + \delta^2\,(|\nabla u|+1)^2 \bigr), \qquad \delta\in(0,1)\\[0.9em]
    &\le \frac12\,\bigl(\frac{1}{\delta^2}\,|Cu|^2 + \delta^2\,(1+\gamma^2)\,|\nabla u|^2 + \delta^2\,(1+\frac{1}{\gamma^2})\bigr)\\[0.9em]
    & = \frac12\,\bigl(\frac{1}{\delta^2}\,|Cu|^2 + 2\delta^2\,|\nabla u|^2 + 2\,\delta^2\bigr) \qquad \text{WLOG setting } \gamma = 1.\\
  \end{align*}
  Returning to the weak form inequality  we get
  \begin{align*}
    \int_\Omega |\nabla u|^2 + \lambda |u|^2\,dx
    &\le \frac12\int_\Omega \frac{f^2}{\varepsilon^2} + \Bigl(\varepsilon^2+\frac{C^2}{\delta^2}\Bigr)u^2
    + 2\delta^2|\nabla u|^2 + 2\delta^2\,dx.\\
  \end{align*}

  Provided $\lambda>\frac{C^2}{2}$, there exist sufficiently small $\varepsilon> 0$ and suitable $\delta < 1$ such that
  \[\lambda > \frac12(\varepsilon^2 + \frac{C^2}{\delta^2}) \qquad\text{ and } \qquad1 > \delta^2\]\\
  hold simultaneously.
  Taking the difference of both sides and denoting the constant factor as $K^2 \coloneqq \int_\Omega \delta^2 \,dx$, yields 
  \begin{align*}
      \alpha^2 \int_\Omega |\nabla u|^2 + |u|^2\,dx
      &\le \frac{1}{2\varepsilon ^2}\int_\Omega f(x,m)^2 \,dx + K^2,
    \end{align*}
  where $\alpha^2$ is the minimum of both factors of $u^2\text{ and } \nabla u^2$. 
Consequently, 
\[\alpha \norm{u}_{\Honez} \le \frac{1}{\varepsilon\sqrt{2}} \norm{f(x,m)}_{\Ltwo}+ K.\]
Thus if $\lambda>\frac{C^2}{2}$, the set 
\[\{u\in \Honez: u=\mu A[u]\ \text{for some } \mu \in[0,1]\}\]
is bounded independently of $\mu$. By \Cref{thm:schaefer}, $A$ has a fixed point $u \in \Honez$ and thus the HJB equation has a solution.
\end{proof}
\begin{remark}
  We can always tune $\varepsilon$ and $\delta$ for given $\lambda>\frac{C^2}{2}$ to satisfy boundedness. This utilises the
  strictness of the inequality in the condition on $\lambda$.
\end{remark}
\begin{proposition}
\label{prop:uniqueness}
If $\lambda>0$ is sufficiently large, then the fixed point $u$ is unique.
\end{proposition}
\begin{proof}
Suppose $u_1,u_2 \in \Honez$ both solve the HJB equation for a given $m \in \Honez$. Taking the weak difference and setting $v=u_1-u_2$ gives
\begin{align*}
  \norm{\nabla(u_1-u_2)}_{L^2}^2+\lambda\norm{u_1-u_2}_{L^2}^2
  &\le \int_\Omega |u_1-u_2|\,|H_1-H_2|\,dx\\[0.7em]
  &\le \tilde{C}\,\norm{u_1-u_2}_{L^2}\,\norm{\nabla(u_1-u_2)}_{L^2}\\
\end{align*}
using Assumption 2 and Cauchy--Schwarz. Hence, by Young's inequality,
\begin{align*}
  \norm{\nabla(u_1-u_2)}_{L^2}^2+\lambda\norm{u_1-u_2}_{L^2}^2 
  &\le \frac12 \bigl(\frac{\tilde{C}^{\,2}}{\varepsilon^2}\,\norm{u_1-u_2}^2_{L^2} +\varepsilon^2 \norm{\nabla(u_1-u_2)}^2_{L^2}\bigr)
  \qquad \varepsilon\in (0, \sqrt{2}).\\
\end{align*}

By the same reasoning as when we proved existence (\Cref{thm:existence_hjb}), 
\[\lambda>\frac{\tilde{C}^{\,2}}{4}\implies u_1=u_2 \text{ in }\Ltwo\iff u_1=u_2 \text{ in }\Honez.\]
\end{proof}
The following theorem summarises our results. 
\begin{theorem}[Existence of unique HJB solution]
  \label{thm:solution_hjb}
  Let the Hamilton-Jacobi equation be as in \cref{eq:system} and suppose that $\lambda > 0$ is sufficiently large. Then for any fixed $m \in \Honez$, there exists
  a unique $u \in \Honez$ solving the equation. 
\end{theorem}
\begin{proof}
  There exists a solution by \Cref{thm:existence_hjb} and it is unique by \Cref{prop:uniqueness}.
\end{proof}
\subsection{Solving the Coupled System of Equations}
Having solved the two equations in (\ref{eq:system}) in isolation, we now look to solve them as a system. As stated above, the fixed point theorems 
will allow us to prove existence of a solution pair for the system as a whole. The challenge is to define and 
prove the existence of a fixed point for a suitable map.
Our work in the previous subsection allows us to define the following maps.

\begin{definition}
\label{def:S-map}
Define $S:\Honez\to\Honez$ by $S[u]=m$, where $m$ solves KFP for given $u$.
\end{definition}
\begin{definition}
\label{def:T-map}
Define $T:\Honez\to\Honez$ by $T[m]=u$, where $u$ solves HJB for given $m$.
\end{definition}
\begin{definition}
\label{def:A-map}
Let $T$ and $S$ be as in \cref{def:S-map} and \cref{def:T-map}. Define
\[A:\Honez\to\Honez,\qquad A \coloneqq S\circ T.\]
\end{definition}

That is, $A$ takes $m\in\Honez$, converts it to $u=T[m]$ solving the HJB equation, and then produces $\widetilde m=S[u]$ solving the KFP equation.
If $A$ has a fixed point, then $m=\widetilde m$, and the coupled system has a solution pair.

We verify the hypotheses of Schaefer's theorem.
\subsubsection{Continuity of the Fixed Point Map}
\begin{lemma}
  \label{thm:A-cont}
  Let $A$ be as in \cref{def:A-map} and suppose that $\lambda > 0$ as it appears in \cref{eq:system} is sufficiently large.
  Then, $A$ is continuous in $\Honez$. 
\end{lemma}
Since $A$ is simply a composition of $S \text{ and } T$, it suffices to show they are both continuous.

\begin{proposition}
  \label{prop:t-cont}
  If $\lambda > 0$ is sufficiently large, $T$ is continuous in $\Honez$. 
\end{proposition}
\begin{proof}
  Let $(m_k)$ converge to $m \in \Honez$. Denote $u=T[m]$, $u_k=T[m_k]$, $H=H(x,\nabla u)$, $H_k = H(x, \nabla u_k)$ and $f_k = f(x, m_k)$.
  Taking the weak formulation of the difference in HJB and setting $v=u-u_k$ gives\\
  \begin{align*}
  \norm{\nabla(u-u_k)}_{L^2}^2+\lambda\norm{u-u_k}_{L^2}^2
  &\le \int_\Omega |u-u_k|\,|H-H_k| + |u-u_k|\,|f-f_k|\,dx\\[0.9em]
  &\le \tilde{C}\norm{u-u_k}_{L^2}\norm{\nabla(u-u_k)}_{L^2}
  + \norm{u-u_k}_{L^2}\norm{f-f_k}_{L^2}. \tag{1}
  \end{align*}
  By Cauchy-Schwarz and Assumption 2. Since $f$ is $\Ltwo$-continuous (Assumption 5) and 
  \[m_k \to m \text{ in } \Honez \implies m_k \to m \text{ in } \Ltwo,\]
  we get $\norm{f-f_k}_{L^2}\to 0$.
  Also, from the energy estimate
  \begin{align*}
  \norm{\nabla u_k}_{L^2}^2+\lambda\norm{u_k}_{L^2}^2
  &\le \int_\Omega  |u_k|\,|f_k| + |u_k|\,|H_k|\,dx.\\
  \end{align*}

Adopting the same argument for boundedness as in \Cref{thm:existence_hjb} gives the inequality
\[\alpha \norm{u_k}_{\Honez} \le \frac{1}{\varepsilon \sqrt{2}} \norm{f_k}_{\Ltwo}+ K\]
 where $K, \,\alpha $ and $\varepsilon$ are independent (positive) constants. By the fact that $f_k \to f$ we know that $\norm{f_k}_{\Ltwo}$ 
 is bounded and so therefore $(u_k)$ is bounded in $\Honez$.
Thus, ${\norm{u-u_k}_{L^2}\norm{f-f_k}_{L^2}\to 0}$. Thus, in the limit of $k$, (1) becomes
\begin{align*}
  \norm{\nabla(u-u_k)}_{L^2}^2+\lambda\norm{u-u_k}_{L^2}^2
  &\le \tilde{C}\norm{u-u_k}_{L^2}\norm{\nabla(u-u_k)}_{L^2}
  \end{align*}
for which $\lambda >0$ being sufficiently large implies $u_k\to u$ i.e. $T$ is continuous in $\Honez$.
\footnote{See the uniqueness argument above (\Cref{prop:uniqueness}).}
\end{proof}
\begin{proposition}
  \label{prop:s-cont}
  If $\lambda > 0$ is sufficiently large, $S$ is continuous in $\Honez$. 
\end{proposition}
\begin{proof}
  Let $(u_k)$ converge to $u$ in $\Honez$. Denote $m=S[u]$, $m_k=S[u_k]$, $b = \Hp(x, \nabla u)$ and $b_k=\Hp(x,\nabla u_k)$.
  Taking the weak formulation of the difference in KFP and setting $v=m-m_k$ yields
  \begin{align*}
  \norm{\nabla(m-m_k)}_{L^2}^2+\lambda\norm{m-m_k}_{L^2}^2
  &\le \int_\Omega |\nabla(m-m_k)|\,\bigl|m\,b - m_k\,b_k\bigr|\,dx.
  \end{align*}

  Define the correction vector field $\beta_k \coloneqq b - b_k$. Note that since $b_k$ is $L^2$ continuous with respect to $u_k$ we have
  $\norm{\beta_k}_{\Ltwo} \to 0$ (Assumption 3). Then, 
  \begin{align*}
  \norm{\nabla(m-m_k)}_{L^2}^2+\lambda\norm{m-m_k}_{L^2}^2
  &\le \int_\Omega |\nabla(m-m_k)|\,\bigl|m\,b - m_k\, b\bigr| +|\nabla(m-m_k)|\,|m_k||\beta_k|\,dx\\
  &\le L \int_\Omega |\nabla(m-m_k)|\,|m - m_k|\,dx + \int_\Omega |\nabla(m-m_k)|\,|m_k|\,|\beta_k|\,dx.
  \end{align*}
  $m$ and $m_k$ are bounded in $\Honez$ by \Cref{thm:lax-milgram} so as $k$ gets large \[\int_\Omega|\nabla(m-m_k)|\,|m_k||\beta_k|\,dx \to 0.\] 
  Therefore, by the largeness of
  $\lambda$ (see \Cref{prop:kfp-coercive}), we get $m_k \to m$ i.e. $S$ is continuous in $\Honez$.
\end{proof}

\begin{remark}
  \label{rm:s-bounded}
  The image of $S$ is always bounded as a direct consequence of Lax--Milgram (\Cref{thm:lax-milgram}). 
\end{remark}
\begin{proof}[Proof of \Cref{thm:A-cont}]
  T is continuous (\Cref{prop:t-cont}) and S is continuous (\Cref{prop:s-cont}) so their composition is continuous in $\Honez$.
\end{proof}
\subsubsection{Compactness of the Fixed Point Map}
\begin{lemma}
  \label{thm:A-comp}
  Let $A$ be as in \cref{def:A-map} and suppose that $\lambda > 0$ as it appears in \cref{eq:system} is sufficiently large.
  Then, $A$ is compact in $\Honez$. 
\end{lemma}
Since $S$ is continuous, it suffices to show that $T$ is compact in $\Honez$:
\[T \text{ compact } \implies A=S\circ T \text{ compact.}\]

We introduce the following theorems. These will allow us to conclude convergence of a subsequence in a bigger space if we are 
given a bounded sequence in a smaller space. The particular relationship that concerns us is the one between $H^2(\Omega) \cap H_0^1(\Omega)$ and $\Honez$.

\begin{definition}
Let $X$ and $Y$ be Banach spaces with $X\subset Y$. We say that $X$ is compactly embedded in $Y$, written $X\subset\subset Y$, provided:
\begin{enumerate}[label=(\roman*), leftmargin=*, itemsep=0.25em]
  \item $\norm{x}_Y \le C\norm{x}_X$ for some constant $C$,
  \item each bounded sequence in $X$ is precompact in $Y$.
\end{enumerate}
\end{definition}

\begin{theorem}[Rellich--Kondrachov compactness theorem]
Assume $U$ is a bounded open subset of $\R^n$ and $\partial U$ is $C^1$. Suppose $1\le p<n$. Then
\[W^{1,p}(U)\subset\subset L^q(U) \quad\text{for each }1\le q<p^\ast, \qquad p^\ast=\frac{np}{n-p}.\]
\end{theorem}

\begin{remark}
Observe $p^\ast>p$. In particular,
\[W^{1,p}(U)\subset\subset L^p(U) \quad\text{for all }1\le p<n.\]
\end{remark}
\begin{corollary}
\label{cor:embed}
\[
H^2(\Omega) \cap H_0^1(\Omega) \subset \subset \Honez.
\]
\end{corollary}
\begin{proof}
Firstly, we have \[\norm{u}_{\Honez} \le \norm{u}_{H^2} \, \text{ for all } u \in H^2(\Omega) \cap \Honez.\]
Now let $(u_k)$ be a bounded sequence in $H^2(\Omega)\cap \Honez$. 
Then $(u_k)$ and $(\partial_i u_k)$ are bounded in $H^1(\Omega)$ for 
each $i=1,\dots,n$. 

By the Rellich--Kondrachov compactness theorem, the embedding 
$H^1(\Omega)\hookrightarrow L^2(\Omega)$ is compact. Hence there exists 
a subsequence, still denoted $(u_k)$, such that
\[
u_k \to u \quad \text{in } L^2(\Omega),
\]
and for each $i$,
\[
\partial_i u_k \to v_i \quad \text{in } L^2(\Omega).
\]

Since weak derivatives are closed under $L^2$ convergence, it follows 
that $v_i=\partial_i u$. Therefore, by a diagonal argument there exists a common subsequence still denoted $(u_k)$ with
\[
u_k \to u \quad \text{strongly in } H^1(\Omega).
\]
Furthermore, $\Honez$ is a closed subspace of $H^1(\Omega)$ so the limit lies in $\Honez$.
\end{proof}
We now have enough to show $T$ is compact.
\begin{proposition}
  \label{prop:t-comp}
  Let $\lambda >0$ be sufficiently large. Then
  \[ |f(x,m)| \le D(|m| + 1) \implies T \text{ is compact in } \Honez.\] 
\end{proposition}
\begin{proof}
  Let $(m_k)$ be a bounded sequence in $\Honez$.
  Denote $u_k = T[m_k]$. We must show that $(u_k)$ admits a subsequence converging strongly in $\Honez$. 
  From \Cref{thm:existence_hjb} we have the inequality
  \[\alpha \norm{u_k}_{\Honez} \le \frac{1}{\varepsilon\sqrt{2}} \norm{f(x,m_k)}_{\Ltwo}+ K\]
  where $\alpha, \varepsilon$ and $K$ are independent (positive) constants. 
  By the growth condition on $f$ we have
  \begin{align*}
    \alpha \norm{u_k}_{\Honez}
    &\le \frac{1}{\varepsilon \sqrt{2}} \norm{f(x,m_k)}_{\Ltwo}+ K\\
    &\le \frac{1}{\varepsilon \sqrt{2}} \bigl(\norm{m_k}_{\Ltwo} + |\Omega|^\frac12\bigr)+ K\\
    &\le \frac{1}{\varepsilon \sqrt{2}} \bigl(\norm{m_k}_{\Honez} + |\Omega|^\frac12\bigr)+ K\\
  \end{align*}
  which is uniformly bounded by our choice of a bounded sequence $(m_k)$ in $\Honez$.
  We also have that 
  \[\norm{H(x,\,\nabla u_k)}_{\Ltwo} \le C(\norm{\nabla u_k}_{\Ltwo} + |\Omega|^\frac12)\] by Assumption 2. This too is uniformly bounded by the boundedness of $(u_k)$ in $\Honez$. This means that both $f_k =f(x,m_k)$
  and $H_k = H(x,\nabla u_k)$ are bounded in $\Ltwo$.
  Recall that $u_k$ weakly solves 
  \[-\Delta u_k +\lambda u_k = g_k,\, g_k \coloneqq f_k - H_k.\]
  By elliptic regularity for the Dirichlet problem, there exists $C>0$ such that 
  \[\norm{u_k}_{H^2(\Omega)} \le C\norm{g_k}_{\Ltwo}\] 
  so $(u_k)$ is bounded in $H^2(\Omega) \cap \Honez$. By \Cref{cor:embed} there exists a subsequence $(u_{k_j})$ converging to $u \in \Honez$. 
  Thus $T$ is compact, and hence $A$ is compact in $\Honez$.
\end{proof}

We have shown that $A$ is continuous and compact in $\Honez$. The last step before we can apply Schaefer's theorem is to prove 
that if there were to exist a fixed point, it would be bounded. 

\subsubsection{Boundedness of the Fixed Point Set}
\begin{lemma}
  \label{thm:A-bounded}
  Let $A$ be as in \cref{def:A-map} then 
  \[\{m\in\Honez: m=\mu A[m]\ \text{for some }\mu\in[0,1]\}\]
  is bounded in $\Honez$.
\end{lemma}
\begin{proof}Consider
\[\{m\in\Honez: m=\mu A[m]\ \text{for some }\mu\in[0,1]\}.\]
The corresponding system, with $A[m]$ substituted for $m/\mu$ \footnote{See \ref{fn:mu}.}, is
\[-\Delta u + H(x,\nabla u) + \lambda u = f(x,m), \qquad \frac{1}{\mu}\Bigl(-\Delta m - \diver\!\bigl(m\,\Hp(x,\nabla u)\bigr) + \lambda m\Bigr)=G(x).\]
Rearranging the latter and using the Lax--Milgram bound implies
\[\norm{m}_{\Honez}\le \frac{1}{c_0}\mu \norm{G}_{L^2}\le \frac{1}{c_0}\norm{G}_{L^2},\]
which is bounded in $\Ltwo$ (Assumption 6). 
\end{proof}
By Schaefer's theorem, $A$ has a fixed point. We may conclude our finding with the proof of \Cref{thm:main_result}.

\paragraph{Proof of the Main Result}
\begin{proof}[Proof of \Cref{thm:main_result}]
Let A be as in \cref{def:A-map}. A is compact and continuous and has bounded fixed points (if they were to exist) by \Cref{thm:A-bounded}, \Cref{thm:A-cont} and \Cref{thm:A-comp}.
By Schaefer's theorem (\Cref{thm:schaefer}), $A$ has a fixed point. Thus, a solution pair $(u,m) \in \Honez \times \Honez$ to the stationary mean field game system exists.
\end{proof}

\subsection{Swapping Composition Order of $T$ and $S$}
In hindsight, we can see that the composition of $T$ (\cref{def:T-map}) followed by $S$ (\cref{def:S-map}) was convenient because it provided a natural bound for the fixed point. We relied on the fact that
the image of $S$ is always bounded (\cref{rm:s-bounded}) independently of $u$.
An equally valid construction for $A$ however, is the reverse composition 
\[A \coloneqq T\circ S.\] 

Here too the existence of a fixed point proves \Cref{thm:main_result} yet crucially, 
recall that we have not shown the image of $T$ to have any constant bound in $\Honez$ rather it depends on the boundedness of its arguments. 
We note that bounded sets under $S$ remain bounded so the new $A$ retains its continuity and compactness.
\footnote{Maps $S$ and $T$ themeselves are unchanged. We have already shown that $S$ is continuous (\Cref{prop:s-cont}) and that $T$ is continuous and compact (\Cref{prop:t-cont,prop:t-comp}).} 
Therefore, we need only demonstrate boundedness of the fixed point set for us to use Schaefer's Theorem. 
Under this new construction consider
\[\{u\in\Honez: u=\mu A[u]\ \text{for some }\mu\in[0,1]\}.\]

The corresponding system, with $A[u]$ substituted for $u/ \mu$ \footnote{See \ref{fn:mu}.}, is
\[-\frac{\Delta u}{\mu} + H(x,\frac{\nabla u}{\mu}) + \frac{\lambda u}{\mu} = f(x,m), \qquad -\Delta m - \diver \bigl(m\,\Hp(x,\nabla u)\bigr) + \lambda m=G(x).\]

We obtain the weak form of HJB with $v = u$
  \begin{align*}
    \int_\Omega |\nabla u|^2 + \lambda |u|^2\,dx
    &= \mu\int_\Omega f u - H(x,\frac{\nabla u}{\mu})u\,dx\\
    &\le \int_\Omega |fu| + \mu\,|H(x,\frac{\nabla u}{\mu})u|\,dx \qquad \text{ since } 0 < \mu \le 1\\
    &\le \int_\Omega |fu| + \mu \,C(|\frac{\nabla u}{\mu}|+1)|u|\,dx \qquad \text{(Assumption 2)}\\
    &= \int_\Omega |fu| + C(|\nabla u|+1)|u|\,dx.
  \end{align*}

This is the same form as the inequality in the proof of \Cref{thm:existence_hjb}. There we arrived at the following bound 

\[\alpha \norm{u}_{\Honez} \le \frac{1}{\varepsilon \sqrt{2}} \norm{f(x,m)}_{\Ltwo}+ K\]
with $\alpha, \, \varepsilon$ and $K$  independent (positive) constants.
The bound of $u$ must be independent of $m$ so we use the growth condition on $f$.
\begin{proposition}
  \label{prop:revA_bounded}
  If $A = T \circ S$, ${\{u\in\Honez: u=\mu A[u]\ \text{for some }\mu\in[0,1]\}}$ is bounded. 
\end{proposition} 
\begin{proof}
  We have
  \begin{align*}
    \alpha \norm{u}_{\Honez}
    &\le \frac{1}{\varepsilon \sqrt{2}} \norm{f(x,m)}_{\Ltwo}+ K\\
    &\le \frac{1}{\varepsilon \sqrt{2}} \bigl(\norm{m}_{\Ltwo} + |\Omega|^\frac12\bigr)+ K\\
    &\le \frac{1}{\varepsilon \sqrt{2}} \bigl(\norm{m}_{\Honez} + |\Omega|^\frac12\bigr)+ K\\
    &\le \frac{1}{\varepsilon \sqrt{2}} \bigl(\frac{1}{c_0}\norm{G}_{\Ltwo} + |\Omega|^\frac12\bigr)+ K
  \end{align*}
  by the Lax--Milgram bound (\Cref{thm:lax-milgram}) since $m \in \operatorname{Im}(S)$. Thus, the set is bounded.
\end{proof}
\begin{remark}
  Notice the similarity in the proofs of \Cref{prop:t-comp} and \Cref{prop:revA_bounded}. They both illustrate that the image of $T$ is bounded on a bounded domain. 
  The difference, when we consider the fixed point of the map as a whole, arises in the domains of $T$. In the first composition $T$ is acting on all $m \in \Honez$ whereas in the latter it acts on 
  the image of $S$. 
\end{remark}
We summarise the analysis of this subsection with a lemma. 
\begin{lemma}
  Let $A= T \circ S$ and suppose that $\lambda > 0$ as it appears in \cref{eq:system} is sufficiently large. 
  Then, $A$ has a fixed point.
\end{lemma}
\begin{proof}
  $A$ is compact and continuous since $T$ is compact and continuous (\Cref{prop:t-comp}, \Cref{prop:t-cont}) and $S$ is bounded continuous (\Cref{prop:s-cont}).
  The set of fixed points is bounded by \Cref{prop:revA_bounded}. 
  Thus, there exists a fixed point by Schaefer's Theorem. 
\end{proof}

\section{Extension to Nonlinear KFP}
In certain applications it is natural to allow for the entry or exit of agents from the system. 
This can be represented by introducing a source term $G(x,m)$ in place of $G(x)$ in \cref{eq:system}. 
Such terms may model, for example, external recruitment of agents into the population, departures from the system, or density–dependent entry and exit effects. 
In this section we consider a class of nonlinear functions $G(x,m)$ satisfying the structural conditions stated in Assumption 7.
Our new system of equations reads
\begin{equation}
\label{eq:system2}
\begin{aligned}
-\Delta u + H(x,\nabla u) + \lambda u &= f(x,m) && \text{in } \Omega,\\
-\Delta m - \diver\!\bigl(m\,\Hp(x,\nabla u)\bigr) + \lambda m &= G(x,m) && \text{in } \Omega.
\end{aligned}
\end{equation}

In this scenario we cannot apply the Lax--Milgram Theorem to the KFP equation. We will therefore adopt fixed point theorems 
to find a solution but as with $H$ and $f$ we impose assumptions on $G(x,m)$.
\begin{enumerate}[label=\arabic*., leftmargin=*, itemsep=0.35em, resume]
\item $G(x,m):\Omega \times \Ltwo \to \Ltwo$ is
  \begin{itemize}[leftmargin=*, itemsep=0.25em]
    \item Lipschitz continuous with respect to $m$ with Lipschitz constant $\tilde{B}$,
    \item bounded by $B\bigl(|m|+1\bigr)$ for a given constant $B$, for all $x\in \Omega$ and $m \in \Ltwo$.
  \end{itemize}
\item In \emph{addition} to our previous constraint on $\lambda$, we have $\lambda > \max(\,B, \,\tilde{B}) +\frac{L^2}{4}$
\end{enumerate}
\begin{example}
For a source term $G$ satisfying Assumption 7, one may consider
\[
G(x,m) = a(x)\,\tanh(m),
\]
where $a \in L^\infty(\Omega)$. From a modelling perspective, this choice corresponds to density-dependent entry 
or exit with saturation: for small $m$ the response is approximately linear, 
while for large $m$ the source term remains bounded.
\end{example}
\paragraph{Originality of the extension.}
To the best of the author's knowledge, the existence result established in this
section for the system with a nonlinear source term $G(x,m)$ does not appear
explicitly in the references cited above. The argument presented here adapts
the fixed-point framework developed for the linear case and shows that the same
method extends under suitable structural assumptions on $G$. The derivation of
the required estimates is carried out explicitly for the present system.

\begin{theorem}[Existence of Nonlinear KFP Solution]
\label{thm:existence2}  
If $\lambda>0$ is sufficiently large, there exists $m\in H^2(\Omega)\cap \Honez$ solving KFP with fixed $u \in \Honez$.
\end{theorem}
\begin{proof}
  We design the map whose fixed point is a solution to the nonlinear equation. 

  \[A:\Honez\to\Honez,\qquad A[m]=w \quad\text{where}\quad -\Delta w -\diver\!\bigl(w\,\Hp(x,\nabla u)\bigr)+\lambda w = G(x,m).\]

  We know $w$ exists by Lax--Milgram \Cref{thm:lax-milgram}, so $A$ is well defined.
  Continuity and compactness are the same as in \Cref{thm:existence_hjb} so we omit them.
  The only difference is in the threshold of $\lambda$ which we'll now derive. As in \Cref{thm:existence_hjb} we want to show
  \[\{m\in \Honez: m=\mu A[m]\ \text{for some } \mu \in [0,1]\}\]
  is bounded.
  Substituting $w=m/\mu$\footnote{See \ref{fn:mu}.} and taking $v=m$, we obtain the weak form
  \begin{align*}
  \int_\Omega |\nabla m|^2 + m \nabla m \Hp + \lambda |m|^2\,dx
  &= \mu\int_\Omega G(x,m)m \,dx\\
  \end{align*}
  Moving the $\Hp$ term over and Assumption 3 give
  \begin{align*}
  \int_\Omega |\nabla m|^2 + \lambda |m|^2\,dx
  &\le \int_\Omega |G(x,m)||m| + L|m||\nabla m|\,dx \qquad \text{since } 0<\mu \le 1\\
  &\le \int_\Omega B(|m|+1)|m| + L|m||\nabla m|\,dx \qquad \text{(Assumption 7)}\\
  \end{align*}
  To isolate $|\nabla m|^2$ and $m^2$ terms on the RHS, we apply Young's inequality \ref{thm:youngs} to give\\
  \begin{align*}
  \int_\Omega |\nabla m|^2 + \lambda |m|^2\,dx
  &\le \int_\Omega \frac12\delta ^2\,|\nabla m|^2 + \bigl(B+\frac12 \varepsilon^2 +\frac{L^2}{2\delta ^2}\bigr)\,|m|^2 \,dx  + K^2
  \qquad \varepsilon >0, \,\delta \in (0, \sqrt{2})\\[0.9em]
  \end{align*}
  where $K^2$ is the sum of terms independent of $m, \nabla m$.
  If $\lambda > B + \frac{L^2}{4}$ then we can find suitable $\varepsilon, \delta$ such that $\norm{m}_{\Honez}$ is bounded. \footnote{See \Cref{thm:existence_hjb} for further elaboration on this point.}
  After taking the difference of both sides, the inequality simplifies to 
  \begin{align*}
    \tilde{\alpha} \norm{m}_{\Honez} \le K
  \end{align*}
  where $\tilde{\alpha}$ is the minimum of both factors of $m^2\text{ and } \nabla m^2$. 
\end{proof}
\begin{remark}
  This result has a very important implication. It tells us that any solution $m \in \Honez$ is still bounded by a constant despite $G(\cdot)$ being nonlinear. 
\end{remark}
\begin{proposition}
If $\lambda>0$ is sufficiently large, then the fixed point $m$ is unique.
\end{proposition}
\label{prop:uniqueness2}
\begin{proof}
Suppose $m_1,m_2$ both solve the nonlinear KFP equation. Taking the difference and setting $v=m_1-m_2$ gives
\begin{align*}
\norm{\nabla(m_1-m_2)}_{L^2}^2 +\lambda\norm{m_1-m_2}_{L^2}^2
&\le \int_\Omega |\Hp|\,|m_1-m_2|\,|\nabla(m_1 -m_2)| + |G_1-G_2|\,|m_1-m_2|\,dx\\
&\le L \norm{m_1-m_2}_{L^2}\norm{\nabla(m_1-m_2)}_{L^2} + \tilde{B}\,\norm{m_1-m_2}_{L^2}^2\\
\end{align*}
using Assumptions 3, 7 and Cauchy--Schwarz. Hence, by our previous use of Young's inequality in \Cref{prop:kfp-coercive},
\[\lambda - \tilde{B} >\frac{L^{\,2}}{4}\implies m_1=m_2 \text{ in }\Honez.\]
\end{proof}

As with the linear version of the KFP equation we define a map $\tilde{S}$ to take the place of $S$ in \cref{def:S-map}.
\begin{definition}
\label{def:s_tilde}
Define $\tilde{S}:\Honez\to\Honez$ by $\tilde{S}(u)=m$, where $m$ solves nonlinear KFP for given $u$.
\end{definition}
Since the form of HJB in \cref{eq:system2} is the same as that in \cref{eq:system}, we can investigate the existence of a fixed point of the map $\tilde{A}$.
\begin{definition}
\label{def:A_tilde}
Let $T$ and $\tilde{S}$ be as in definitions \ref{def:S-map} and \ref{def:T-map}. Define
\[\tilde{A}:\Honez\to\Honez,\qquad A \coloneqq \tilde{S}\circ T.\]
\end{definition}

A fixed point of $\tilde{A}$ means a solution to the new system of equations exists. 
\paragraph{Main Result for Nonlinear KFP}
\begin{theorem}[Existence of MFG Solution with Nonlinear KFP]
Let the Mean Field Game system be characterised by \cref{eq:system2} and suppose the assumptions above are true. 
Then a solution $(u,m) \in \Honez \times \Honez$ exists. 
\end{theorem}
\begin{proof}
  We prove the theorem via showing $\tilde{A}$ has a fixed point. 
  Once again, we wish to use Schaefer's theorem. We need to show that $\tilde{S}$ is continuous. 
  Let $(u_k)$ converge to $u \in \Honez$. Denote $m=\tilde{S}[u]$, $m_k=\tilde{S}[u_k]$, $b=\Hp(x, \nabla u)$, $b_k=\Hp(x,\nabla u_k)$ and $G_k = G(x,m_k)$.\\
  Taking the weak formulation of the difference in KFP and setting $v=m-m_k$ yields
  \begin{align*}
  \norm{\nabla(m-m_k)}_{L^2}^2+\lambda\norm{m-m_k}_{L^2}^2
  &\le \int_\Omega |\nabla(m-m_k)|\,\bigl|m\,b - m_k\,b_k\bigr| +\, |G - G_k|\,|m-m_k|\,dx.
  \end{align*}
  Applying the Lipschitz condition (Assumption 7) to $G,G_k$ gives
  \begin{align*}
  \norm{\nabla(m-m_k)}_{L^2}^2+\lambda\norm{m-m_k}_{L^2}^2
  &\le \int_\Omega |\nabla(m-m_k)|\,\bigl|m\,b - m_k\,b_k\bigr| +\, \tilde{B}|m-m_k|^2\,dx.
  \end{align*}

  By $\lambda$ being sufficiently large to account for the $\tilde{B}$ term (Assumption 8) and recalling the same line of reasoning in \Cref{prop:s-cont} we conclude that 
  $\tilde{S}$ is continuous in $\Honez$. 
  We also have that $\tilde{S}$ is bounded (\Cref{thm:existence2}) and that $T$ is both compact in $\Honez$ (\Cref{prop:t-comp}) and continuous in $\Honez$ (\Cref{prop:t-cont}). $\tilde{A}$ is therefore
  compact in $\Honez$.\\

  To meet the final condition of Schaefer's theorem, we must show 
  \[\{m\in\Honez: m=\mu \tilde{A}[m]\ \text{for some }\mu\in[0,1]\}\]
  is bounded. 
  The corresponding system, with $\tilde{A}[m]$ substituted for $m/\mu$ \footnote{See \ref{fn:mu}.}, is
  \[-\Delta u + H(x,\nabla u) + \lambda u = f(x,m), \qquad \frac{1}{\mu}\Bigl(-\Delta m - \diver\!\bigl(m\,\Hp(x,\nabla u)\bigr) + \lambda m\Bigr)=G(x,\,\frac{m}{\mu}).\]
  
  Taking the weak form of the latter and setting $v=m$ implies
  \begin{align*}
  \int_\Omega |\nabla m|^2 + \lambda |m|^2\,dx
  &\le \int_\Omega \mu|G(x,\,\frac{m}{\mu})||m| + L|m||\nabla m|\,dx \qquad \text{See \Cref{thm:existence2}}\\
  &\le \int_\Omega B(|m|+1)|m| + L|m||\nabla m|\,dx \qquad \text{(Assumption 7)}\\
  \end{align*}
  In \Cref{thm:existence2} we showed that when $\lambda>0$ is sufficiently large, the inequality simplifies to 
  \[\tilde{\alpha} \norm{m}_{\Honez} \le K\]
  Thus, the set 
  \[\{m\in\Honez: m=\mu \tilde{A}[m]\ \text{for some }\mu\in[0,1]\}\]
  is bounded independently of $\mu$. 
  By Schaefer's theorem (\ref{thm:schaefer}), $\tilde{A}$ has a fixed point.
\end{proof}

\section{Conclusion}
We briefly summarise the strategy used to establish the existence of steady state solutions to the coupled system. 
Our approach proceeds by analysing the two equations separately and then combining the resulting solution operators. 
For a fixed density $m \in \Honez$, the Hamilton--Jacobi--Bellman equation is solved to obtain $u \in \Honez$, 
while for a fixed $u \in \Honez$ the Kolmogorov--Fokker--Planck equation determines the corresponding density $m \in \Honez$. 
These steps define two operators between the relevant function spaces. 
By composing these operators we obtain a map whose fixed points correspond to solution pairs $(u,m)$ of the coupled system.\\

The fixed point framework therefore provides a flexible mechanism for 
passing from solvability of the individual equations to the existence of a solution for the full mean field game system. 
Finally, we further extend this analysis by incorporating nonlinear source terms $G(x,m)$ in the KFP equation and show that, 
under suitable structural assumptions, the same argument yields existence of steady state solutions in this more general setting.

\bibliographystyle{plain}
\bibliography{references}
\end{document}